\documentclass{amsart}
% For dealing with references we use the comment environment
\usepackage{verbatim}
\newenvironment{reference}{\comment}{\endcomment}
%\newenvironment{reference}{}{}
\newenvironment{slogan}{\comment}{\endcomment}
\newenvironment{history}{\comment}{\endcomment}

% For commutative diagrams we use Xy-pic
\usepackage[all]{xy}

% We use 2cell for 2-commutative diagrams.
\xyoption{2cell}
\UseAllTwocells

% We use multicol for the list of chapters between chapters
\usepackage{multicol}

% This is generally recommended for better output
\usepackage{lmodern}
\usepackage[T1]{fontenc}

% For cross-file-references

% Package for hypertext links:
\usepackage{hyperref}

% For any local file, say "hello.tex" you want to link to please

% Theorem environments.
%
\theoremstyle{plain}
\newtheorem{theorem}[subsection]{Theorem}
\newtheorem{proposition}[subsection]{Proposition}
\newtheorem{lemma}[subsection]{Lemma}

\theoremstyle{definition}
\newtheorem{definition}[subsection]{Definition}
\newtheorem{example}[subsection]{Example}
\newtheorem{exercise}[subsection]{Exercise}
\newtheorem{situation}[subsection]{Situation}

\theoremstyle{remark}
\newtheorem{remark}[subsection]{Remark}
\newtheorem{remarks}[subsection]{Remarks}

\numberwithin{equation}{subsection}

% Macros
%
\def\lim{\mathop{\mathrm{lim}}\nolimits}
\def\colim{\mathop{\mathrm{colim}}\nolimits}
\def\Spec{\mathop{\mathrm{Spec}}}
\def\Hom{\mathop{\mathrm{Hom}}\nolimits}
\def\Ext{\mathop{\mathrm{Ext}}\nolimits}
\def\SheafHom{\mathop{\mathcal{H}\!\mathit{om}}\nolimits}
\def\SheafExt{\mathop{\mathcal{E}\!\mathit{xt}}\nolimits}
\def\Sch{\mathit{Sch}}
\def\Mor{\mathop{\mathrm{Mor}}\nolimits}
\def\Ob{\mathop{\mathrm{Ob}}\nolimits}
\def\Sh{\mathop{\mathit{Sh}}\nolimits}
\def\NL{\mathop{N\!L}\nolimits}
\def\CH{\mathop{\mathrm{CH}}\nolimits}
\def\proetale{{pro\text{-}\acute{e}tale}}
\def\etale{{\acute{e}tale}}
\def\QCoh{\mathit{QCoh}}
\def\Ker{\mathop{\mathrm{Ker}}}
\def\Im{\mathop{\mathrm{Im}}}
\def\Coker{\mathop{\mathrm{Coker}}}
\def\Coim{\mathop{\mathrm{Coim}}}

% Boxtimes
%
\DeclareMathSymbol{\boxtimes}{\mathbin}{AMSa}{"02}

%
% Macros for moduli stacks/spaces
%
\def\QCohstack{\mathcal{QC}\!\mathit{oh}}
\def\Cohstack{\mathcal{C}\!\mathit{oh}}
\def\Spacesstack{\mathcal{S}\!\mathit{paces}}
\def\Quotfunctor{\mathrm{Quot}}
\def\Hilbfunctor{\mathrm{Hilb}}
\def\Curvesstack{\mathcal{C}\!\mathit{urves}}
\def\Polarizedstack{\mathcal{P}\!\mathit{olarized}}
\def\Complexesstack{\mathcal{C}\!\mathit{omplexes}}
% \Pic is the operator that assigns to X its picard group, usage \Pic(X)
% \Picardstack_{X/B} denotes the Picard stack of X over B
% \Picardfunctor_{X/B} denotes the Picard functor of X over B
\def\Pic{\mathop{\mathrm{Pic}}\nolimits}
\def\Picardstack{\mathcal{P}\!\mathit{ic}}
\def\Picardfunctor{\mathrm{Pic}}
\def\Deformationcategory{\mathcal{D}\!\mathit{ef}}

\setlength{\emergencystretch}{3em}
\begin{document}
\title[Stacks Project, Tag 091V]{411 --- Construction of derived completion as a reflective subcategory}
\maketitle
\noindent Copyright (C) 2005--2025 Johan de Jong. Stacks Project authors. Source: \href{https://stacks.math.columbia.edu/tag/091V}{Tag 091V}.

\noindent Editorial difficulty note (not author text). Difficulty H2: The author constructs completion by derived Hom from an extended Cech complex. Localization annihilates that complex in the relevant directions, proving derived completeness, and a triangle argument verifies the full universal mapping property rather than just defining a candidate object..

\noindent Editorial provenance note (not author text). History: excerpt from revision \texttt{a0444\allowbreak 6e57e\allowbreak c1fbc\allowbreak 252a8\allowbreak 71afc\allowbreak ec775\allowbreak 2fb28\allowbreak 07b14}, more-algebra.tex, lines 26153--26231. Reference presentation was adapted; original-author context and core support blocks were retained or added. The mathematical wording is unchanged. Licensed under GNU FDL 1.2 or later; no invariant sections or cover texts. See ../licenses/COPYING. Context blocks reproduce author definitions and supporting results; they are not additional counted problems.

\begin{lemma}
\label{lemma-derived-completion}
\begin{slogan}
Derived completions along finitely generated ideals exist, and can
be computed by a {\v C}ech procedure.
\end{slogan}
Let $I$ be a finitely generated ideal of a ring $A$.
The inclusion functor $D_{comp}(A, I) \to D(A)$ has a
left adjoint, i.e., given any object $K$ of $D(A)$ there
exists a map $K \to K^\wedge$ of $K$ into a derived complete
object of $D(A)$ such that the map
$$
\Hom_{D(A)}(K^\wedge, E) \longrightarrow \Hom_{D(A)}(K, E)
$$
is bijective whenever $E$ is a derived complete object of $D(A)$.
In fact, if $I$ is generated by $f_1, \ldots, f_r \in A$, then we have
$$
K^\wedge = R\Hom\left((A \to \prod\nolimits_{i_0} A_{f_{i_0}} \to
\prod\nolimits_{i_0 < i_1} A_{f_{i_0}f_{i_1}}
\to \ldots \to A_{f_1\ldots f_r}), K\right)
$$
functorially in $K$.
\end{lemma}

\begin{proof}
Define $K^\wedge$ by the last displayed formula of the lemma.
There is a map of complexes
$$
(A \to \prod\nolimits_{i_0} A_{f_{i_0}} \to
\prod\nolimits_{i_0 < i_1} A_{f_{i_0}f_{i_1}} \to
\ldots \to A_{f_1\ldots f_r}) \longrightarrow A
$$
which induces a map $K \to K^\wedge$. It suffices to prove that
$K^\wedge$ is derived complete and that $K \to K^\wedge$ is an
isomorphism if $K$ is derived complete\footnote{Namely, if $E$
is derived complete and $a : K \to E$ is a map, then the
commutative diagram
$$
\xymatrix{
K \ar[r]_a \ar[d] & E \ar[d]^{\cong} \\
K^\wedge \ar[r]^{a^\wedge} & E^\wedge
}
$$
shows that any map into $E$ factors through $K^\wedge$.
Choose a distinguished triangle $K \to K^\wedge \to C$ and apply
${}^\wedge$. We find that $C^\wedge = 0$.
By the above this means that $\Hom(C, E) = 0$ for $E$ derived complete
and this proves that $\Hom(K, E) = \Hom(K^\wedge, E)$ as desired.}.

\medskip\noindent
Let $f \in A$. By Lemma \href{https://stacks.math.columbia.edu/tag/0A6C}{lemma double localize (Tag 0A6C)}
the object $R\Hom_A(A_f, K^\wedge)$ is equal to
$$
R\Hom\left((A_f \to \prod\nolimits_{i_0} A_{ff_{i_0}} \to
\prod\nolimits_{i_0 < i_1} A_{ff_{i_0}f_{i_1}} \to
\ldots \to A_{ff_1\ldots f_r}), K\right)
$$
If $f \in I$, then $f_1, \ldots, f_r$ generate the
unit ideal in $A_f$, hence the extended alternating
{\v C}ech complex
$$
A_f \to \prod\nolimits_{i_0} A_{ff_{i_0}} \to
\prod\nolimits_{i_0 < i_1} A_{ff_{i_0}f_{i_1}} \to
\ldots \to A_{ff_1\ldots f_r}
$$
is zero in $D(A)$ by
Lemma \href{https://stacks.math.columbia.edu/tag/0G6K}{lemma extended alternating torsion (Tag 0G6K)}.
(In fact, if $f = f_i$ for some $i$, then this complex
is homotopic to zero by
Lemma \href{https://stacks.math.columbia.edu/tag/0G6J}{lemma extended alternating homotopy zero (Tag 0G6J)};
this is the only case we need.)
Hence $R\Hom_A(A_f, K^\wedge) = 0$ and we conclude that
$K^\wedge$ is derived complete by Lemma \href{https://stacks.math.columbia.edu/tag/091P}{lemma hom from Af (Tag 091P)}.

\medskip\noindent
Conversely, if $K$ is derived complete, then $R\Hom_A(A_f, K)$
is zero for all $f = f_{i_0} \ldots f_{i_p}$, $p \geq 0$. Thus
$K \to K^\wedge$ is an isomorphism in $D(A)$.
\end{proof}
\end{document}
